\documentclass[11pt,a4paper]{article}
\usepackage[utf8]{inputenc}
\usepackage[T1]{fontenc}
\IfFileExists{spanish.ldf}{%
  \usepackage[spanish,es-nodecimaldot]{babel}
}{%
  \usepackage{babel}
}
\usepackage{amsmath,amssymb,bm}
\usepackage[margin=2.4cm]{geometry}
\usepackage{booktabs,longtable,array}
\usepackage[hidelinks]{hyperref}
\setlength{\emergencystretch}{3em}
\newcommand{\dd}{\mathrm{d}}
\newcommand{\Id}{\mathbf{I}}
\title{Modelo de movimiento colectivo celular\\Unidades y guía de elección de parámetros}
\author{Documento de trabajo}
\date{}
\begin{document}
\maketitle

\section{Objetivo}

En este documento vamos a mostrar los parámetros que debemos elegir para las simulaciones.

\section{Geometría, superposición y vecinos}

Lo primero que podemos pensar es en qué momento consideramos que dos células interactúan. Para eso, recordemos primero algunas definiciones del modelo.

Todas las células conservan un área $a=\pi R^2$, donde $R$ es el radio de una célula redonda. Para una relación de aspecto $\phi_i\geq1$, los semiejes son

\begin{equation}
 \ell_{\parallel,i}=R\sqrt{\phi_i},\qquad
 \ell_{\perp,i}=\frac{R}{\sqrt{\phi_i}}.
\end{equation}

La posición y orientación son $\bm r_i$ y $\varphi_i$. Definimos

\begin{equation}
 \alpha_i=\ell_{\parallel,i}^{2}+\ell_{\perp,i}^{2},\qquad
 \epsilon_i=\frac{\phi_i^2-1}{\phi_i^2+1},\qquad
 \bm Q_i=\begin{pmatrix}
 \cos2\varphi_i&\sin2\varphi_i\\\sin2\varphi_i&-\cos2\varphi_i
 \end{pmatrix}.
\end{equation}

Para $\bm r_{ij}=\bm r_i-\bm r_j$:

\begin{equation}
 \bm M_{ij}=\frac{\alpha_i\epsilon_i\bm Q_i+\alpha_j\epsilon_j\bm Q_j}{\alpha_i+\alpha_j},
 \qquad
 \beta_{ij}=(\alpha_i+\alpha_j)^2-(\alpha_i\epsilon_i-\alpha_j\epsilon_j)^2
 -4\alpha_i \alpha_j\epsilon_i\epsilon_j\cos^2(\varphi_i-\varphi_j).
\end{equation}

Introducimos la matriz:
\begin{equation}
 \bm G_{ij}=\frac{\alpha_i+\alpha_j}{\beta_{ij}}(\Id-\bm M_{ij}).
\end{equation}

La superposición entre dos células es entonces:
\begin{equation}
	I_{ij}=I_{ij}^{0}\xi_{ij}
\end{equation}

Donde $I_{ij}^{0}$ es el máximo para esas formas y orientaciones,y $\xi_{ij}$ la superposición normalizada:

\begin{equation}
 I_{ij}^{0}=\frac{4a^2}{\pi\sqrt{\beta_{ij}}},\qquad
 \xi_{ij}=\frac{I_{ij}}{I_{ij}^{0}}
 =\exp(-\bm r_{ij}^{\mathsf T}\bm G_{ij}\bm r_{ij}).
\end{equation}

Definimos que dos células interactúan cuando la superposición normalizada excede un umbral:

\begin{equation}
 \boxed{\xi_{ij}>\xi_{\mathrm{cut}}.}
\end{equation}

Por lo tanto, debemos elegir $\xi_{\mathrm{cut}}$. La justificación de la elección se desarrolla en el análisis de vecinos.

Para buscar los vecinos, nos limitamos a buscar en la grilla espacial que selecciona candidatos antes de evaluar $\xi_{ij}$. Su alcance debe cubrir todas las parejas que puedan superar el umbral, basta un alcance

\begin{equation}
 r_{\mathrm{cut}}^{\max}
 =2R\sqrt{\phi_{\max}[-\ln(\xi_{\mathrm{cut}})]}.
\end{equation}

La implementación periódica
requiere al menos tres celdas de grilla por eje.

\section{Interacción, fuerza y torque}

Definimos la energía de interacción como:
\begin{equation}
 u_\lambda(\xi)=\kappa\left[(1-\lambda)\xi^\gamma
 -\lambda\ln(1-\xi^\gamma)\right],\qquad 0\leq\lambda\leq1.
\end{equation}

Obtenemos la fuerza y el torque de la siguiente forma:
\begin{equation}
 \bm f_{ij}=-\nabla_{\bm r_{ij}}u_\lambda,\qquad
 \tau_{ij}=-\partial_{\varphi_i}u_\lambda.
\end{equation}
Definiendo
\begin{equation}
 B_\lambda(\xi)=1+\lambda\frac{\xi^\gamma}{1-\xi^\gamma},
\end{equation}
resulta
\begin{equation}
 \bm f_{ij}=2\kappa\gamma\xi_{ij}^{\gamma}B_\lambda(\xi_{ij})
 \bm G_{ij}\bm r_{ij},
\end{equation}
\begin{equation}
 \tau_{ij}=\kappa\gamma\xi_{ij}^{\gamma}B_\lambda(\xi_{ij})\bm r_{ij}^{\mathsf T} \frac{\partial\bm G_{ij}}{\partial\varphi_i}\bm r_{ij}.
\end{equation}

El origen de la regularización y su comportamiento a distancias cortas se
discuten en el documento de fuerzas.

\section{Movimiento y elección de unidades}

La dinámica es sobreamortiguada. Definimos los vectores
\begin{equation}
	\bm e_i=(\cos\varphi_i,\sin\varphi_i),
	\qquad
	\bm e_i^\perp=(-\sin\varphi_i,\cos\varphi_i),
\end{equation}
que indican las direcciones paralela y perpendicular al eje
de autopropulsión de la célula.

La parte determinista de las ecuaciones de movimiento es
\begin{equation}
	\frac{\mathrm{d}\bm r_i}{\mathrm{d}t}
	=
	v_e h(\phi_i)\bm e_i
	+
	\bm\mu_i
	\sum_{j\in\Omega_i}\bm f_{ij},
	\label{eq:dimensional_translation}
\end{equation}
\begin{equation}
	\frac{\mathrm{d}\varphi_i}{\mathrm{d}t}
	=
	\mu_{r,i}
	\sum_{j\in\Omega_i}\tau_{ij},
	\label{eq:dimensional_rotation}
\end{equation}
donde
\begin{equation}
	h(\phi_i)=\frac{\phi_i-1}{\phi_{\max}-1}.
\end{equation}
En el modelo instantáneo, $h=0$ para las células redondas
y $h=1$ para las elongadas. Por tanto, $v_e$ es la velocidad
de autopropulsión de una célula completamente elongada,
no necesariamente su velocidad total.

\subsection{Escalas de movilidad e interacción}

Sea $\mu_0$ la movilidad traslacional de referencia de una
célula redonda. Separamos la escala dimensional de movilidad
de su dependencia con la forma:
\begin{equation}
	\bm\mu_i=\mu_0\bm m_i,
	\qquad
	\mu_{r,i}=\frac{\mu_0}{R^2}m_r(\phi_i),
\end{equation}
con
\begin{equation}
	\bm m_i
	=
	m_\parallel(\phi_i)\bm e_i\bm e_i^{\mathsf T}
	+
	m_\perp(\phi_i)\bm e_i^\perp
	(\bm e_i^\perp)^{\mathsf T}.
\end{equation}
Las funciones $m_\parallel$, $m_\perp$ y $m_r$ son
adimensionales.

La energía de interacción puede escribirse como
\begin{equation}
	u_{ij}=\kappa\,U_{ij},
\end{equation}
donde $\kappa$ tiene unidades de energía y $U_{ij}$ es una función adimensional de la geometría relativa, $\gamma$ y $\lambda$.

Introducimos las posiciones adimensionales
\begin{equation}
	\widetilde{\bm r}_i=\frac{\bm r_i}{R}.
\end{equation}
Como
\begin{equation}
	\nabla_{\bm r_i}
	=
	\frac{1}{R}\nabla_{\widetilde{\bm r}_i},
\end{equation}
la fuerza y el torque toman la forma
\begin{equation}
	\bm f_{ij}
	=
	\frac{\kappa}{R}\bm F_{ij},
	\qquad
	\bm F_{ij}
	=
	-\nabla_{\widetilde{\bm r}_i}U_{ij},
\end{equation}
\begin{equation}
	\tau_{ij}=\kappa\,T_{ij},
	\qquad
	T_{ij}=-\frac{\partial U_{ij}}{\partial\varphi_i}.
\end{equation}
La fuerza contiene un factor $1/R$ por la derivada espacial.
El torque no lo contiene porque el ángulo es adimensional.

\subsection{Reescalado de las ecuaciones de movimiento}

Elegimos como unidad temporal el tiempo que una célula
con velocidad de autopropulsión $v_e>0$ tarda en recorrer
una distancia $R$:
\begin{equation}
	t_0=\frac{R}{v_e},
	\qquad
	\widetilde t=\frac{t}{t_0}=\frac{v_e}{R}t.
\end{equation}
Entonces,
\begin{equation}
	\frac{\mathrm{d}\bm r_i}{\mathrm{d}t}
	=
	v_e
	\frac{\mathrm{d}\widetilde{\bm r}_i}
	{\mathrm{d}\widetilde t},
	\qquad
	\frac{\mathrm{d}\varphi_i}{\mathrm{d}t}
	=
	\frac{v_e}{R}
	\frac{\mathrm{d}\varphi_i}{\mathrm{d}\widetilde t}.
\end{equation}

Sustituyendo las escalas de fuerza y movilidad en
la ecuación traslacional, obtenemos
\begin{equation}
	v_e
	\frac{\mathrm{d}\widetilde{\bm r}_i}
	{\mathrm{d}\widetilde t}
	=
	v_e h(\phi_i)\bm e_i
	+
	\frac{\mu_0 \kappa}{R}
	\bm m_i
	\sum_{j\in\Omega_i}\bm F_{ij}.
\end{equation}
Dividiendo por $v_e$,
\begin{equation}
	\frac{\mathrm{d}\widetilde{\bm r}_i}
	{\mathrm{d}\widetilde t}
	=
	h(\phi_i)\bm e_i
	+
	\frac{\mu_0 \kappa}{Rv_e}
	\bm m_i
	\sum_{j\in\Omega_i}\bm F_{ij}.
	\label{eq:scaled_translation}
\end{equation}

Para la orientación, el mismo procedimiento da
\begin{equation}
	\frac{v_e}{R}
	\frac{\mathrm{d}\varphi_i}{\mathrm{d}\widetilde t}
	=
	\frac{\mu_0 \kappa}{R^2}
	m_r(\phi_i)
	\sum_{j\in\Omega_i}T_{ij}.
\end{equation}
Multiplicando por $R/v_e$,
\begin{equation}
	\frac{\mathrm{d}\varphi_i}{\mathrm{d}\widetilde t}
	=
	\frac{\mu_0 \kappa}{Rv_e}
	m_r(\phi_i)
	\sum_{j\in\Omega_i}T_{ij}.
	\label{eq:scaled_rotation}
\end{equation}

Por tanto, ambas ecuaciones contienen el mismo
acoplamiento adimensional:
\begin{equation}
	\boxed{
		\widetilde\kappa=\frac{\mu_0 \kappa}{Rv_e}
	}.
	\label{eq:kappa_definition}
\end{equation}

Las ecuaciones deterministas quedan
\begin{equation}
	\boxed{
		\frac{\mathrm{d}\widetilde{\bm r}_i}
		{\mathrm{d}\widetilde t}
		=
		h(\phi_i)\bm e_i
		+
		\widetilde\kappa\bm m_i
		\sum_{j\in\Omega_i}\bm F_{ij}
	},
\end{equation}
\begin{equation}
	\boxed{
		\frac{\mathrm{d}\varphi_i}{\mathrm{d}\widetilde t}
		=
		\widetilde\kappa m_r(\phi_i)
		\sum_{j\in\Omega_i}T_{ij}
	}.
\end{equation}

Por lo tanto, en el código podemos fijar el radio de las células redondas y la velocidad de autopropulsión de las elongadas a uno:

\begin{equation}
	\widetilde R=1,
	\qquad
	\widetilde v_e=1.
\end{equation}

Estos dejan de aparecer como un parámetro independiente de las ecuaciones deterministas: su efecto relativo a la interacción queda incorporado en $\widetilde\kappa$.

A $R$ y $\mu_0$ fijos, aumentar $\kappa$ incrementa $\widetilde\kappa$, mientras que aumentar $v_e$ lo reduce. Así, $\widetilde\kappa$ compara la escala de velocidad inducida por las interacciones, $\mu_0\kappa/R$, con la velocidad de autopropulsión.

\subsection{Relación con los parámetros del código}

Entonces si en las simulaciones elegimos $\widetilde R=\widetilde v_e=1$, el parámetro \texttt{kRep} representa numéricamente $\widetilde \kappa$.

\subsection{Reescalado del ruido}

Para explicitar también el cambio de unidades del ruido,
consideramos coeficientes dimensionales de difusión
\begin{equation}
	D_{\parallel,i}=D_0m_\parallel(\phi_i),
	\qquad
	D_{\perp,i}=D_0m_\perp(\phi_i),
	\qquad
	D_{r,i}=D_{r,0}m_r(\phi_i).
\end{equation}
Los incrementos aleatorios durante un intervalo
$\Delta t$ se escriben como
\begin{equation}
	\begin{split}
		\Delta\bm r_{i,\mathrm{ruido}}
		={}&
		\sqrt{2D_0m_\parallel\Delta t}\,
		Z_\parallel\bm e_i\\
		&+
		\sqrt{2D_0m_\perp\Delta t}\,
		Z_\perp\bm e_i^\perp,
	\end{split}
\end{equation}
\begin{equation}
	\Delta\varphi_{i,\mathrm{ruido}}
	=
	\sqrt{2D_{r,0}m_r\Delta t}\,Z_r,
\end{equation}
donde las variables $Z_\parallel$, $Z_\perp$ y $Z_r$
son normales estándar independientes en cada paso
y para cada célula.

Usando
\begin{equation}
	\Delta t=\frac{R}{v_e}\Delta\widetilde t,
	\qquad
	\Delta\widetilde{\bm r}=\frac{\Delta\bm r}{R},
\end{equation}
obtenemos
\begin{equation}
	\begin{split}
		\Delta\widetilde{\bm r}_{i,\mathrm{ruido}}
		={}&
		\sqrt{\frac{2D_0}{Rv_e}}\,
		\sqrt{m_\parallel\Delta\widetilde t}\,
		Z_\parallel\bm e_i\\
		&+
		\sqrt{\frac{2D_0}{Rv_e}}\,
		\sqrt{m_\perp\Delta\widetilde t}\,
		Z_\perp\bm e_i^\perp,
	\end{split}
\end{equation}
\begin{equation}
	\Delta\varphi_{i,\mathrm{ruido}}
	=
	\sqrt{
		2\frac{D_{r,0}R}{v_e}
		m_r\Delta\widetilde t
	}\,Z_r.
\end{equation}

Definimos entonces los parámetros de ruido del código:
\begin{equation}
	\boxed{
		\eta^2=\frac{2D_0}{Rv_e},
		\qquad
		d_\varphi=\frac{D_{r,0}R}{v_e}
	}.
\end{equation}
Aquí $\eta$ corresponde a \texttt{noise\_eta} y
$d_\varphi$ a \texttt{d\_phi}. La difusividad angular
adimensional efectiva es $d_\varphi m_r$.

Así, en presencia de ruido, la dependencia con $v_e$
queda incorporada no solamente en $\kappa$, sino
también en los parámetros adimensionales del ruido.
Cambiar $v_e$ manteniendo fijos los coeficientes
dimensionales de difusión modifica esos parámetros.

\subsection{Paso temporal y período de deformación}

También debemos estudiar el paso temporal de integración y un paso temporal que nos diga cada cuánto tiempo se intentan deformar las células. Expresándo esto en las mismas unidades:

\begin{equation}
	\Delta\widetilde t=\frac{v_e}{R}\Delta t,
	\qquad
	\widetilde T_{\mathrm{def}}
	=\frac{v_e}{R}T_{\mathrm{def}}.
\end{equation}

\subsection{Conclusión}

Los parámetros relativos a la fuerza y las ecuaciones de movimiento que debemos estudiar para ver cómo varía el sistema son:
\begin{equation}
	\widetilde \kappa, \qquad \gamma, \qquad \lambda, \qquad \eta, \qquad d_\varphi, \qquad \Delta \widetilde t, \qquad \widetilde T_{def} 
\end{equation}

Mientras que quedan fijados:

\begin{equation}
	\widetilde R=1,
	\qquad
	\widetilde v_e=1.
\end{equation}

\section{Reglas de deformación}
En cada barrido de deformación, una célula redonda prueba $n_{\mathrm{int}}$
orientaciones. Cada propuesta parte del mismo estado original, desplaza
el centro en la dirección propuesta una distancia igual al aumento del
semieje mayor y conserva el área. Se admite si
\begin{equation}
 \max_{j\in\mathcal C_p}\xi_{ij}^{(p)}\leq\xi_{\mathrm{cut}},
\end{equation}

donde $\mathcal C_p$ es el conjunto de candidatos de la grilla para la propuesta. Entre las propuestas válidas se elige la de menor máximo; los empates se sortean. Por tanto, aumentar $n_{\mathrm{int}}$ puede modificar tanto la probabilidad de elongación como la posición seleccionada.

La contracción se solicita cuando la velocidad determinista proyectada sobre el eje de autopropulsión no es positiva. Se propone mantener el centro y reducir la relación de aspecto. Si se activa la condición de seguridad, se rechaza cuando el máximo normalizado de la configuración propuesta supera $\xi_{\mathrm{safe}}$.

\section{Estado de las decisiones}
\small
\begin{longtable}{p{0.20\textwidth}p{0.26\textwidth}p{0.44\textwidth}}
\toprule
Cantidad & Estado o valor & Justificación y trabajo pendiente \\
\midrule
\endhead
$\widetilde R$, $\widetilde v_e$ & $1$, $1$ & Fijos por el reescaleo. \\
$\phi_{\max}$ & $2.7$ & Motivación experimental. \\
$\gamma$ & $3$ & Valor usado en papers anteriores. \\
$\xi_{\mathrm{cut}}$ & $e^{-1}$ & Concuerda mucho con la distancia de interacción, desarrollado en el análisis de vecinos. \\
$\lambda$ & $0.1$ & Candidato motivado por el umbral teórico de monotonía, desarrollado en el análisis de fuerzas. \\
$\xi_{\mathrm{safe}}$ & $0.8$ & Cercano al punto problemático, desarrollado en el análisis de fuerzas. \\
$\kappa$ & Por elegir & . \\
$\eta$, $d_\varphi$ & Por documentar/elegir & . \\
Condición inicial & Aleatoria, redondas; $t_{\mathrm{stab}}=0$ & Protocolo principal. Sus transitorios forman parte del análisis desde el tiempo inicial. \\
Deformación & Instantánea & Elegido como un primer modelo. \\
$n_{\mathrm{int}}$ & Por elegir & . \\
Período de deformación & Por fijar & . \\
$\Delta t$ & $0.05$, exploratorio & Validar convergencia para los valores de interacción seleccionados; no asumir estabilidad para cualquier $\kappa$. \\
\bottomrule
\end{longtable}
\normalsize

\end{document}
